% 从现有最终论文提取的个人版式；与题目内容无关的样式集中在此。
% 原正式版本中的按页 Needspace、图片裁剪等局部补丁未复制。
% 字体优先保持原稿在macOS上的宋体/黑体；无此字体时使用TeX发行版Fandol。
\IfFontExistsTF{Songti SC}{
  \setCJKmainfont{Songti SC}[BoldFont=Songti SC Bold]
  \setCJKsansfont{Heiti SC}
  \setCJKmonofont{Songti SC}
  \setCJKfamilyfont{zhsong}{Songti SC}
  \setCJKfamilyfont{zhhei}{Heiti SC}
  \setCJKfamilyfont{zhkai}{Songti SC}
  \setCJKfamilyfont{zhfs}{Songti SC}
}{
  \setCJKmainfont{FandolSong-Regular.otf}[BoldFont=FandolSong-Bold.otf]
  \setCJKsansfont{FandolHei-Regular.otf}[BoldFont=FandolHei-Bold.otf]
  \setCJKmonofont{FandolSong-Regular.otf}
  \setCJKfamilyfont{zhsong}{FandolSong-Regular.otf}
  \setCJKfamilyfont{zhhei}{FandolHei-Regular.otf}
  \setCJKfamilyfont{zhkai}{FandolKai-Regular.otf}
  \setCJKfamilyfont{zhfs}{FandolFang-Regular.otf}
}
\providecommand{\songti}{\CJKfamily{zhsong}}
\providecommand{\heiti}{\CJKfamily{zhhei}}
\providecommand{\kaishu}{\CJKfamily{zhkai}}
\providecommand{\fangsong}{\CJKfamily{zhfs}}
\usepackage{needspace}
\usepackage{fvextra}
\usepackage{placeins}
\usepackage{microtype}
\DeclareMathSizes{9.5}{9.5}{8}{8}
\newcommand{\important}[1]{{\heiti\bfseries #1}}
\newcommand{\questionlabel}[1]{{\heiti\bfseries\fontsize{14}{18}\selectfont #1}}
\newcommand{\topic}[1]{{\heiti\bfseries\fontsize{11}{14}\selectfont #1}}
\newcommand{\questionline}[1]{\par\noindent\hspace*{2em}{\heiti\bfseries\fontsize{14}{18}\selectfont #1\par}}
\setmainfont{texgyretermes-regular.otf}[BoldFont=texgyretermes-bold.otf,ItalicFont=texgyretermes-italic.otf,BoldItalicFont=texgyretermes-bolditalic.otf]
\setsansfont{texgyreheros-regular.otf}[BoldFont=texgyreheros-bold.otf,ItalicFont=texgyreheros-italic.otf,BoldItalicFont=texgyreheros-bolditalic.otf]
\IfFontExistsTF{Menlo}{\setmonofont{Menlo}}{\setmonofont{lmmono10-regular.otf}}
\xeCJKDeclareCharClass{CJK}{"22A5}
\geometry{a4paper,margin=25mm}
\renewcommand{\baselinestretch}{1.28}
\renewcommand{\thesection}{\chinese{section}、}
\renewcommand{\thesubsection}{\arabic{section}.\arabic{subsection}}
\renewcommand{\thefigure}{1-\arabic{figure}}
\renewcommand{\thetable}{1-\arabic{table}}
\makeatletter
\renewcommand{\mcm@cap@keywordsname}{关键词}
\renewcommand\keywords[1]{%
  \renewcommand{\mcm@tokens@keywords}{#1}%
  \par\vskip1ex
  {\noindent\zihao{-4}\heiti\bfseries\mcm@cap@keywordsname：}~{\mcm@tokens@keywords}%
}
\renewcommand\section{\@startsection{section}{1}{\z@}%
 {-3.0ex \@plus -1ex \@minus -.2ex}{1.8ex \@plus .2ex}%
 {\centering\normalfont\heiti\fontsize{14}{18}\selectfont\bfseries}}
\renewcommand\subsection{\@startsection{subsection}{2}{\z@}%
 {-1.4ex \@plus -.5ex \@minus -.2ex}{.5ex \@plus .1ex}%
 {\normalfont\heiti\fontsize{14}{18}\selectfont\bfseries}}
\renewcommand\subsubsection{\@startsection{subsubsection}{3}{2em}%
 {-1.5ex \@plus -.5ex \@minus -.2ex}{.7ex \@plus .2ex}%
 {\normalfont\heiti\fontsize{11}{14}\selectfont\bfseries}}
\makeatother
\captionsetup[figure]{font={small,normalfont},labelfont={bf},textfont={normalfont},skip=5pt}
\captionsetup[table]{font={small,normalfont},labelfont={bf},textfont={normalfont},skip=5pt}
\widowpenalty=10000 \clubpenalty=10000 \displaywidowpenalty=10000
\makeatletter
\setlength{\@fptop}{0pt}\setlength{\@fpsep}{14pt}\setlength{\@fpbot}{0pt plus 1fil}
\makeatother
\setlength{\textfloatsep}{12pt plus 3pt minus 2pt}
\setlength{\intextsep}{12pt plus 3pt minus 2pt}
\setlength{\emergencystretch}{2em}
\pagestyle{plain}
% 每问图表采用“问号-序号”；总路线图延续原稿计入第一问图序。
\newcommand{\startquestion}[3]{%
  \FloatBarrier
  \subsection{问题#2}
  \renewcommand{\thefigure}{#1-\arabic{figure}}%
  \renewcommand{\thetable}{#1-\arabic{table}}%
  \ifnum#1>1\setcounter{figure}{0}\fi
  \setcounter{table}{0}%
  \questionline{第#2问：#3}%
}
% 仅用于模板说明，不自动进入未来正式论文。
\newcommand{\templatehint}[1]{\noindent{\small\textcolor{gray}{［模板说明］#1}}\par}

% 保留短题名单行下划线；长题名自然换行，避免不可断行下划线盒溢出。
\newlength{\personalTitleWidth}
\newcommand{\papertitlecontent}{%
  \settowidth{\personalTitleWidth}{\PaperTitle}%
  \ifdim\personalTitleWidth>\linewidth
    \PaperTitle\par\nobreak\vskip3pt\hrule
  \else
    \underline{\PaperTitle}%
  \fi
}
